\chapter[Band 7: Surjektive Funktionen]{Band 7 auf einen Blick}
\noindent\textbf{Surjektive Funktionen}\par\medskip
\noindent\emph{Lesart.} Surjektivität ist ein eigenes Strukturaxiom
für \(F\colon A\to B\); sie hängt ausdrücklich von der Zielmenge \(B\) ab.
\begin{BandUebersicht}
\textbf{Surjektive Funktion}
\(F\colon A\sur B\) bedeutet \(F\colon A\to B\) zusammen mit
\UebFormel{y\in B\vdash\exists x\in A\;F(x)=y.}
\UebQuellen{\FormulaRefAuto{Surjektive Funktion}[def];
\FormulaRefAuto{Surjektivität}[ax]}
&
\textit{Surjektivität ist gleichbedeutend mit vollem Bild}
Für \(F\colon A\to B\):
\UebFormel{\begin{aligned}
&F\colon A\sur B\leftrightarrow F[A]=B,\\
&F\colon A\sur F[A].
\end{aligned}}
\UebQuellen{\FormulaRefAuto{SurjectiveIffImageEqualsCodomain}[thm];
\FormulaRefAuto{CorestrictionToImageSurjective}[thm]}\\
\addlinespace[.8ex]
\textbf{Potenzmenge und Diagonalmenge}
Für \(F\colon A\to\powerset(A)\) betrachtet der Beweis
\UebFormel{D\coloneqq\{x\in A\mid x\notin F(x)\}.}
Dies ist eine Aussonderung in \(A\), keine zusätzliche Axiomannahme.
&
\textit{Cantors Satz}
\UebFormel{\begin{aligned}
&F\colon A\to\powerset(A)\\
&\hspace{10mm}\vdash\neg(F\colon A\sur\powerset(A)).
\end{aligned}}
Keine Funktion auf einer Menge trifft alle ihre Teilmengen.
\UebQuellen{\FormulaRefAuto{CantorNoPowerSetSurjection}[thm]}\\
\addlinespace[.8ex]
\textbf{Surjektion aus einer Injektion}
Für \(F\colon A\inj B\), \(a_0\in A\) wird
\(\Gsurjfrominj{F}{a_0}\) durch Rückzuordnung auf \(F[A]\)
und den festen Wert \(a_0\) außerhalb von \(F[A]\) definiert:
\UebFormel{\begin{aligned}
&(b,a)\in\Gsurjfrominj{F}{a_0}\ \leftrightarrow\\
&b\in B\land a\in A\land\bigl(\\
&\quad(b\in F[A]\land F(a)=b)\\
&\quad\lor(b\notin F[A]\land a=a_0)\bigr).
\end{aligned}}
\UebQuellen{\FormulaRefAuto{F\colon A\inj B \vdash \Gsurjfrominj{F}{a_0} \coloneqq \{\, (b,a)\in B\times A \mid (b\in F[A]\land F(a)=b)\lor (b\notin F[A]\land a=a_0)\,\}}[def]}
&
\textit{Die Rückzuordnung ist surjektiv}
\UebFormel{\begin{aligned}
&F\colon A\inj B\dsep a_0\in A\\
&\hspace{10mm}\vdash\Gsurjfrominj{F}{a_0}\colon B\sur A.
\end{aligned}}
Die ausdrücklich verlangte Wahl von \(a_0\) behandelt die Werte
außerhalb des Bildes.
\UebQuellen{\FormulaRefAuto{F\colon A\inj B \dsep a_0\in A \vdash \Gsurjfrominj{F}{a_0}\colon B\sur A}[thm]}\\
\addlinespace[.8ex]
\textbf{Projektionen}
Für \((x,y)\in A\times B\):
\UebFormel{\pi_1(x,y)\coloneqq x,\qquad
\pi_2(x,y)\coloneqq y.}
\UebQuellen{\FormulaRefAuto{FirstProjectionDef}[def];
\FormulaRefAuto{SecondProjectionDef}[def]}
&
\textit{Surjektivität bei nichtleerem anderem Faktor}
\UebFormel{\begin{aligned}
&b_0\in B\vdash\pi_1\colon A\times B\sur A,\\
&a_0\in A\vdash\pi_2\colon A\times B\sur B.
\end{aligned}}
\UebQuellen{\FormulaRefAuto{FirstProjectionSurjective}[thm];
\FormulaRefAuto{SecondProjectionSurjective}[thm]}\\
\addlinespace[.8ex]
\textbf{Fasern und Einschränkung auf eine Relation}
\UebFormel{\Fib_F(y)\coloneqq F^{-1}[\{y\}].}
Für eine totale Relation \(\TotRel{R,A,B}\) wird außerdem
\(\pi_1\) auf \(R\) eingeschränkt.
\UebQuellen{\FormulaRefAuto{FiberMapDef}[def];
\FormulaRefAuto{RestrictionDef}[def]}
&
\textit{Nichtleere Fasern und surjektive Relationsprojektion}
\UebFormel{\begin{aligned}
&F\colon A\sur B\dsep y\in B\vdash\Fib_F(y)\neq\varnothing,\\
&\TotRel{R,A,B}\vdash\pi_1\restriction_R\colon R\sur A.
\end{aligned}}
\UebQuellen{\FormulaRefAuto{F\colon A\sur B \dsep y\in B \vdash \Fib_F(y)\neq\varnothing}[thm];
\FormulaRefAuto{\pi_1\restriction_R\colon R\sur A}[thm]}\\
\addlinespace[.8ex]
\textbf{Komposition}
Für \(F\colon A\sur B\), \(G\colon B\sur C\):
\UebFormel{(G\circ F)(x)\coloneqq G(F(x))\quad(x\in A).}
\UebQuellen{\FormulaRefAuto{CompositionDef}[def]}
&
\textit{Surjektivität bleibt erhalten}
\UebFormel{G\circ F\colon A\sur C.}
Auch der leere Randfall ist ausgearbeitet:
\UebFormel{\varnothing\colon\varnothing\sur\varnothing.}
\UebQuellen{\FormulaRefAuto{G\circ F\colon A\sur C}[thm];
\FormulaRefAuto{\varnothing\colon \varnothing\sur \varnothing}[thm]}\\
\end{BandUebersicht}
